\chapter{Présentation du projet}

\section{Contexte}

Table complète est une application Android développée en Kotlin et Jetpack
Compose dans le cadre du M1 à l'ESGI.

L'idée de départ vient de Theo : une application autour des cocktails,
construite sur l'API publique TheCocktailDB. En cherchant des sources de
données, nous avons trouvé TheMealDB, son équivalent pour la cuisine. Les deux
bases se complètent naturellement, l'une donne le plat, l'autre le verre :
nous avons donc décidé de les réunir.

Le sujet exclut tout backend. Les données viennent d'API publiques et de
fichiers embarqués dans l'application, et toute la logique (filtres, accords,
menus, liste de courses) s'exécute sur le téléphone. Ce cadre nous a orientés
vers un projet où l'intérêt tient surtout à l'interface et à la manière de
manipuler ces données.

\section{Objectifs}

Nous voulons une application ludique et interactive, qui donne envie de
cuisiner plutôt qu'un simple moteur de recherche de recettes. Elle doit :

\begin{itemize}
  \item proposer chaque jour un repas complet : un plat, un verre et un dessert ;
  \item adapter les suggestions au régime, au matériel et aux boissons de
        l'utilisateur ;
  \item accompagner la préparation étape par étape, jusqu'à la liste de courses ;
  \item présenter les recettes comme une carte de restaurant, avec des
        animations et des retours haptiques soignés.
\end{itemize}

Côté projet, l'objectif est une base de code propre, testée et construite à
deux sur Git.

\section{Public cible}

L'application s'adresse à toute personne qui cuisine chez elle et cherche une
idée de repas. Aucun niveau n'est requis : chaque recette se suit pas à pas.

Les accords avec du vin ou des cocktails alcoolisés sont réservés aux
personnes majeures. Les autres utilisateurs ont accès à toute l'application,
avec uniquement des boissons sans alcool.

\section{Équipe}

Le projet est réalisé par deux étudiants et encadré par M. Ramzy Kermad.
Sans backend, tout le travail porte sur l'application. Nous répartissons donc
la charge par fonctionnalité, de l'écran jusqu'aux données et aux tests, pour
que chacun intervienne sur toutes les couches.

\begin{tableau}{L{3.4cm} Y}
  \entete{Membre} & \entete{Périmètre principal} \\
  \midrule
  Mamadou Cisse & Plats : exploration, fiche recette, mode cuisine, liste de courses \\
  \midrule
  Theo Robert & Boissons : cocktails, vins, accords, composition de menus \\
  \midrule
  En commun & Architecture, design system, premier lancement, favoris, hors ligne, relecture des PR \\
\end{tableau}
